\documentclass[11pt, twocolumn, landscape, a4paper]{article}

\usepackage[margin=1cm]{geometry}
\usepackage[utf8]{inputenc}
%\usepackage[spanish]{babel}
\usepackage[spanish,es-noshorthands]{babel}
\usepackage{amssymb, amsmath, amsbsy}
\usepackage{tasks}
\usepackage{enumerate}
\usepackage{enumitem}
\usepackage{graphicx}

\usepackage{tikz}
\usetikzlibrary{angles,quotes}
%\usetikzlibrary {arrows.meta} 
%\usetikzlibrary{fit,positioning}
\usepackage{multicol}
\usepackage[pdftex]{hyperref}
\hypersetup{pdftitle={Razonamiento matemático},pdfauthor={Luis Carrillo Gutiérrez}}

\fontfamily{cmss}\selectfont
\renewcommand{\familydefault}{\sfdefault}
\pagestyle{empty}

\begin{document}
% Evaluación del aprendizaje
\subsection*{Razonamiento matemático}
\begin{itemize}
\item{Analizar \iffalse y responde \fi que en el reloj análogo/analógico de la abuela son las 03:00 p.m. ¿Cuál es la medida del ángulo (interno) que describen las manecillas en ese instante?}
\item{Estimar, nombrar y clasifícar la medida de cada uno de los siguientes ángulos. Luego medir y verificar la estimación.

\begin{multicols}{2}
{\scriptsize Figura a.}

%\shorthandoff{"}
\begin{tikzpicture}
	\coordinate (A) at (-1.75, -3.1);
	\coordinate (B) at (0, 0);
	\coordinate (C) at (3.5, 0);
	
	\begin{scope}[-latex]
	\draw (B) -- (C);
	\draw (B) -- (A);
	\end{scope}	
	\draw pic[draw, angle radius=.35cm] {angle=A--B--C};
\end{tikzpicture}


{\scriptsize Figura b.}

\begin{tikzpicture}
	\coordinate (A) at (1.75, -3.1);
	\coordinate (B) at (0, 0);
	\coordinate (C) at (3.5, -.75);
	
	\begin{scope}[-latex]
	\draw (B) -- (C);
	\draw (B) -- (A);
	\end{scope}	
	\draw pic[draw, angle radius=.6cm] {angle=A--B--C};
\end{tikzpicture}
\end{multicols}
}
%\item{Rosa hace la siguiente afirmación.}
\item{Si dos recta paralelas son cortadas simultáneamente por una recta transversal, donde se forman ocho ángulos.

{\scriptsize Figura c.}

\begin{tikzpicture}
	\begin{scope}[latex-latex]
	\draw (-2, 0) -- (4, 0);
	\draw (-2, -1.5) -- (4, -1.5);
	\draw (-1, 1.25) -- (3, -2.25);
	\end{scope}
	\draw (-1.2, 1.45) node{$A$};
	\draw (3.2, -2.45) node{B};
	\draw (-1.9, .275) node{$W$};
	\draw (3.9, .275) node{$Z$};
	\draw (-1.9, -1.175) node{$V$};
	\draw (4, -1.175) node{$T$};
	
	\draw (.75, .275) node{$X$};
	\draw (1.975, -1.85) node{$Y$};
\end{tikzpicture}

\begin{itemize}
\item[$\circ$]{¿Encontrar/identificar parejas de ángulos congruentes en la figura anterior que son (y no) opuestos por el vértice? Explicar/usar un transportador.}

\end{itemize}
}
\end{itemize}
\end{document}